
The following definitions are presented from standard mathematical sources (\textbf{Sets}).
Clarifications and labels have been added where applicable.

\begin{define}{Basic Set Concepts}\phantomsection\label{def:basic-set-concepts} % \hyperref[def:basic-set-concepts]{YOUR_CONTEXT}
  \begin{enumerate}
    \item \textbf{Universal Set:} A universal set $\setUNIV$ is a
    set that contains all objects under consideration in a particular context or problem.
    \item \textbf{Empty Set:} The empty set, denoted $\setNULL$,
    is the unique set that contains no elements.
    \item \textbf{Subset:} Let $A$ and $B$ be sets. $A$ is a
    subset of $B$ ($A \subseteq B$) if (and only if)
    for all $x \in A$, $x \in B$.
    \item \textbf{Set Equality:} Let $A$ and $B$ be sets. The sets $A$ and $B$ are
    equal ($A = B$) if $A \subseteq B$ and $B \subseteq A$.
    \item \textbf{Power Set:} Let $A$ be a set. The power set of $A$, denoted $\sPOWER{A}$ or $2^A$,
    is the set of all subsets of $A$. That is, $\mathcal{P}(A) = \{S\in\setUNIV : S \subseteq A\}$.
  \end{enumerate}
\end{define}
The following definitions extend the understanding of set operations
and properties.


\begin{define}{Set Operations}\phantomsection\label{def:set-operations} % \hyperref[def:set-operations]{YOUR_CONTEXT}
  \begin{enumerate}
    \item \textbf{Complementary Set:} Let $A$ be a subset of the universal set $\setUNIV$.
    The complement of $A$, denoted $A^c$ or $\overline{A}$, is the set of all elements
    in $\setUNIV$ that are not in $A$. That is, $A^c = \{x \in \setUNIV : x \notin A\}$.
    \item \textbf{Union:} Let $A$ and $B$ be sets. The union of $A$ and $B$,
    denoted $A \cup B$, is the set of all elements that are in $A$, or in $B$, or in both.
    That is, $A \cup B = \{x : x \in A \lor x \in B\}$.
    \item \textbf{Intersection:} Let $A$ and $B$ be sets. The intersection of $A$ and $B$,
    denoted $A \cap B$, is the set of all elements that are in both $A$ and $B$.
    That is, $A \cap B = \{x : x \in A \land x \in B\}$.
    \item \textbf{Set Difference:} Let $A$, $B$ be subsets of a universal set $\setUNIV$.
    The set difference $A \sMINUS B$ is defined to be the set
    $\{ x \in \setUNIV \LDIV x \in A \text{ and } x \notin B\}$. \vspace*{10px}\\
    \textit{It is sometimes called the relative complement of $B$ in $A$,
    and denoted in the textbook as $A - B$.}
    \item \textbf{Cartesian Product:} Let $A$ and $B$ be sets. The Cartesian product of
    $A$ and $B$, denoted $A \sTIMES B$, is the set consisting of all ordered pairs:
    \[
    A \sTIMES B = \{(a, b) \LDIV a \in A, b \in B\}.
    \]
  \end{enumerate}
\end{define}


\newpage

The following claims are given in the reference as well. Let $A,B\in U$, a universal set.


\begin{claim}{Subset of Union: $A \subseteq \sOR{A}{B}$}\phantomsection\label{claim:subset-union} % \hyperref[claim:subset-union]{YOUR_CONTEXT}
  \begin{subprove}[bluecolor]
      Let $x\in U$, assume $x\in A$. Then the statement \inlinequote{$x\in A$ or $x\in B$} is true.
      Hence, by the definition, $\sOR{A}{B}=\{x\in U\LDIV x\in A$ or $x\in B\}$, $x\in \sOR{A}{B}$.
      Thus, $A \subseteq \sOR{A}{B}$.
  \end{subprove}
\end{claim}

\begin{claim}{Empty Subset: $\setNULL \subseteq A$}\phantomsection\label{claim:empty-subset} % \hyperref[claim:empty-subset]{YOUR_CONTEXT}
  \begin{subprove}[bluecolor]
      Let $x\in U$, assume $x\in \setNULL$. But there are no elements in $\setNULL$, so
      $x\in \setNULL$ is false. Since the premise is false, the implication if $x\in \setNULL$,
      then $x\in A$ is true. Thus, $\setNULL \subseteq A$ by definition of \hyperref[def:basic-set-concepts]{subsets}.
  \end{subprove}
\end{claim}

\begin{claim}{Subset as Equality: $\sOR{A}{B}=A$ if and only if $B\subseteq A$}\phantomsection\label{claim:subset-equality} % \hyperref[claim:subset-equality]{YOUR_CONTEXT}
  \begin{subprove}[bluecolor]
      Let $x\in U$, first, we assume $B\subseteq A$, we have already shown
      $A \subseteq \sOR{A}{B}$ by the claim, \hyperref[claim:subset-union]{Subset of Union}. 
      To show $\sOR{A}{B} \subseteq A$, assume $x\in \sOR{A}{B}$.
      By definition, $x\in A$ or $x\in B$. Considering $x\in B$ and assumption
      $B \subseteq A$, $x\in B$ implies $x\in A$. Thus, $\sOR{A}{B} \subseteq A$.
      Hence, if $B\subseteq A$, then $\sOR{A}{B} = A$. \vspace*{10px}\\
      Now, we assume $\sOR{A}{B}=A$. To show $B\subseteq A$, we now assume $x\in B$. Since $x\in B$, by definition
      of union, $x\in \sOR{A}{B} = A$. Therefore, $B\subseteq A$.
  \end{subprove}
\end{claim}


%%%%%%%%%%%%%%%%%%% 9/8/25

Let $\setUNIV$ be a set. 
\begin{claim}{ICA6}\phantomsection\label{claim:ica6} % \hyperref[claim:ica6]{YOUR_CONTEXT}
  For all $A,B\in\setUNIV$, $$A\sTIMES B =\setNULL\THEN \group{A=\setNULL\OR B=\setNULL}$$
  \begin{subprove}[bluecolor]
    Let $A,B\in\setUNIV$. We prove $A\sTIMES B=\setNULL$. Assume it is not the case that 
    $A=\setNULL$ or $B=\setNULL$. Thus, we have $A=\setNULL$ and $B=\setNULL$. 
    Since $A\ne\setNULL$ by assumption, there exists an element $a\in A$. Likewise, there 
    is some $b\in B$ since $B\ne\setNULL$. Thus, $\group{a,b}\in\sBUILD{(x,y)}{x\in A,y\in B}$ 
    by definition. But then, $A\sTIMES B\ne\setNULL$ by definition. Hence, the claim holds. 
  \end{subprove}
\end{claim}


\begin{define}{Relatively Prime Integers (Coprime)}\phantomsection\label{def:relatively-prime} % \hyperref[def:relatively-prime]{YOUR_CONTEXT}
  Let $x,y \in \mathbb{Z}$. The integers $x$ and $y$ are said to be
  \textbf{relatively prime} if and only if
  $\gcd(x,y) = 1$.
\end{define}
